\section{本章小结}
\label{sec:dens-summary}

本章开头问：一个测度何时能写成参考测度的加权积分，权重怎样计算，有了权重之后积分怎样改写。全章可以对照有限点集上的向量来读，表~\ref{tab:dens-vector-analogy}列出了对应关系，以及类比失效之后一般情形用什么代替。

\begin{BookInlineTable}
\small
\begin{tabular}{@{}>{\RaggedRight\arraybackslash}p{22mm}>{\RaggedRight\arraybackslash}p{\dimexpr(\linewidth-22mm-4\tabcolsep)*2/5\relax}>{\RaggedRight\arraybackslash}p{\dimexpr(\linewidth-22mm-4\tabcolsep)*3/5\relax}@{}}
\toprule
概念 & 有限点集上的向量 & 一般测度\\
\midrule
测度 & 非负向量 & 集合上的可数可加函数\\
绝对连续 & 零分量位置的包含 & 对每个可测集检查零测关系\\
相互奇异 & 两向量无公共正分量 & 用一个可测集分开，闭支撑可以相交\\
密度 & 逐分量比值 & 无逐点比值，由集合上的积分等式确定\\
Lebesgue分解 & 按$m_i>0$与否分下标 & 寻找承载最多质量的零测集\\
换参考 & 对角矩阵乘向量 & 换测度公式\\
连续换参考 & 对角矩阵相乘 & 链式规则，几乎处处成立\\
对调参考与目标 & 对角矩阵求逆 & 倒数规则，要求两个测度等价\\
\bottomrule
\end{tabular}
\caption{密度理论的向量类比。}
\label{tab:dens-vector-analogy}
\end{BookInlineTable}

回答从零测集开始。绝对连续$\nu\ll\mu$要求参考测度忽略的集合目标测度也忽略，它有方向，而且必须对每个可测集检查，单点检查在不可数集合上不够。相互奇异$\mu\perp\nu$要求一个可测集把两种质量完全分开，闭支撑相交的两个测度仍可以相互奇异。

加权方向总是可行的：任给非负可测权重，对参考测度积分就得到一个测度，并且它相对于参考绝对连续。反方向由Radon–Nikodym定理回答：两个测度都$\sigma$有限时，$\nu\ll\mu$保证密度存在，密度在$\mu$几乎处处意义下唯一。证明取所有不超过目标的权重中积分最大的一个，Hahn分解说明剩余质量只能为零。计数测度的反例说明参考测度的$\sigma$有限性不能随意去掉；有限有符号测度则对两个变差分别求密度，得到一个可积密度，其正负部对应Jordan分解。

一般的测度未必绝对连续。Lebesgue分解把它唯一地拆成绝对连续部分与奇异部分：构造依靠承载最多质量的零测集，唯一性只读取两个集合上的值，密度表示再用Radon–Nikodym定理。实轴上以长度为参考时，奇异部分又分成原子与奇异连续部分。累积函数的跳跃读出原子，Cantor测度说明连续的累积函数仍可以把全部质量放在长度为零的集合上，普通导数只恢复绝对连续部分的密度。

有了密度，换测度公式把$\int f\,\mathrm d\nu$改写成$\int fh\,\mathrm d\mu$：点和被积函数都不变，只是累计权重乘上$h$。相对同一参考的密度可以非负组合，带符号的组合需要有限性。逐次换参考时密度相乘，各因子的零测例外最终都是底层参考的零测例外。两个测度等价时密度可以取倒数，单向绝对连续时只能在正密度区域上取倒数，零密度区域上的参考质量无法恢复。

具体计算时先查绝对连续，再求密度。离散测度的密度是参考质量为正处的点质量之比，参考零测点上若仍有目标质量，密度不存在，应先分出奇异部分。实轴上单调累积函数绝对连续时，导数就是密度；一般情形先读出跳跃，再检查去掉跳跃后的部分是否绝对连续。本章的所有运算都在同一个空间上进行；下一章研究质量沿可测映射从一个空间转移到另一个空间，那时目标空间上的密度要结合本章的换测度公式与坐标换元来计算。
